%% macros_finale.tex — Macro condivise per la tesi finale (Zamar, 2026-07-20)
%% Le date e i numeri canonici fanno riferimento alla run overnight 2026-05-13
%% (resume 2026-05-14 post-reboot), single source of truth: STATE.md root.

\providecommand{\cleanroomrun}{run overnight 2026-05-13\xspace}

\providecommand{\engineversion}{3.4.0\xspace}
\providecommand{\schemaversion}{3.3\xspace}
\providecommand{\llmmodel}{Llama~3.1~8B~Q4\_K\_M\xspace}
%% Nei nomi lunghi dei modelli si inseriscono punti di rottura (\allowbreak)
%% dopo i trattini/slash: i token \texttt non spezzabili causavano overfull.
\providecommand{\embedmodel}{\texttt{paraphrase-\allowbreak multilingual-\allowbreak mpnet-\allowbreak base-\allowbreak v2}\xspace}
\providecommand{\reranker}{\texttt{nickprock/\allowbreak cross-\allowbreak encoder-\allowbreak italian-\allowbreak bert-\allowbreak stsb}\xspace}
\providecommand{\nlimodel}{mDeBERTa-\allowbreak v3-\allowbreak base-\allowbreak xnli-\allowbreak multilingual-\allowbreak nli-\allowbreak 2mil7\xspace}

%% Conteggi canonici.
\providecommand{\goldn}{122\xspace}
\providecommand{\goldRimb}{69\xspace}
\providecommand{\goldNonRimb}{39\xspace}
\providecommand{\goldNonDet}{9\xspace}
\providecommand{\goldRouted}{5\xspace}

%% Identificatori delle quattro Note AIFA codificate nel sistema:
%%   N01 = inibitori di pompa protonica (IPP) per gastroprotezione cronica
%%   N13 = statine nella dislipidemia (ipolipemizzanti)
%%   N66 = antinfiammatori non steroidei (FANS)
%%   N97 = anticoagulanti orali diretti (DOAC) nella fibrillazione atriale
%%         non valvolare (FANV)
\providecommand{\notezero}{N01}
\providecommand{\notetredici}{N13}
\providecommand{\notesessanta}{N66}
\providecommand{\notenova}{N97}
%% Conteggi test verificati con pytest il 2026-07-08 (435 + 590 = 1025).
\providecommand{\testcountengine}{435\xspace}
\providecommand{\testcountorch}{590\xspace}
\providecommand{\testcountall}{1025\xspace}

%% Sezioni con piccole maiuscole per le decisioni cliniche.
\providecommand{\decision}[1]{\textsc{#1}}

%% Comandi per pesi ALES.
\providecommand{\alesalpha}{0{,}40}
\providecommand{\alesbeta}{0{,}30}
\providecommand{\alesgamma}{0{,}20}
\providecommand{\alesdelta}{0{,}10}
